\documentclass[11pt,a4paper]{article}
\usepackage[margin=2.2cm]{geometry}
\usepackage{amsmath,amssymb,bm}
\usepackage{booktabs,array,longtable}
\usepackage{enumitem}
\usepackage{microtype}
\usepackage{xcolor}
\usepackage[colorlinks,linkcolor=red!50!black,citecolor=red!70!black,urlcolor=red!80!black]{hyperref}

\newcommand{\bbH}{\mathbb{H}}
\newcommand{\bbT}{\mathbb{T}}
\newcommand{\ket}[1]{|#1\rangle}
\newcommand{\bra}[1]{\langle#1|}
\newcommand{\braket}[1]{\langle#1\rangle}
\newcommand{\code}[1]{\texttt{#1}}
\newcommand{\vR}{v^{R}}
\newcommand{\vbR}{\bar v^{R}}
\newcommand{\so}{\mathrm{so}}
\newcommand{\act}{\mathrm{act}}

\title{Equations implemented in \texttt{numerics/mrgw\_hermitian}\\[2mm]
\large A reimplementation guide for the Hermitian pole-space MR-$GW$ code}
\author{}
\date{September 22, 2026}

\begin{document}
\maketitle

\section*{Purpose and reading guide}

This document writes out, in the order in which the code evaluates them, all
equations implemented in the package \code{numerics/mrgw\_hermitian} and in
the three scripts \code{run\_checks.py}, \code{example\_h4.py} and
\code{check\_ozone.py}.  Every formula is stated with the index conventions,
normalizations, thresholds and sign choices actually used, so that the
methods can be reimplemented without reading the code.  The theory behind the
equations is in \code{notes/mr-hedin.tex} (Objectives~I and~II),
\code{notes/comparison.tex} and Wang, Fang, and Li, arXiv:2604.16013 (WFL
below); equation numbers of the form (S$n$) refer to the Supplemental
Material of WFL.  Section~\ref{sec:index} maps each equation to the function
that implements it.

Throughout, all quantities are real (real orbitals, real CI vectors), and
$\dagger$ can be read as a transpose.  Energies are in Hartree inside the
code; the scripts convert to eV with $1\,\mathrm{Ha}=27.211386245988$~eV.

\section{Conventions and PySCF primitives}
\label{sec:conv}

\paragraph{Spin orbitals.}
Spatial molecular orbitals $p_s=0,\ldots,n_{\rm mo}-1$ are ordered core,
active, virtual (indices $0\le p_s<n_{\rm core}$, $n_{\rm core}\le p_s<n_{\rm occ}$,
$p_s\ge n_{\rm occ}=n_{\rm core}+n_{\rm cas}$).  Spin orbitals are labelled
\begin{equation}
  p=2p_s+\sigma,\qquad \sigma=0\ (\alpha),\ 1\ (\beta),\qquad n_{\so}=2n_{\rm mo}.
  \label{eq:so}
\end{equation}
The index sets $\mathcal C$ (core), $\mathcal A$ (active), $\mathcal V$
(virtual) and $\mathcal I=\mathcal C\cup\mathcal V$ (inactive) are the
corresponding lists of spin orbitals; ``local active'' indices
$x_\ell=2(x_s-n_{\rm core})+\sigma\in\{0,\ldots,n_{\act}-1\}$,
$n_{\act}=2n_{\rm cas}$, label the active spin orbitals inside the CAS
routines.  Symbols: $P,Q$ inactive; $i,j$ core; $a,b$ virtual; $x,y,z,w$
active; $p,q,r,s$ arbitrary.

\paragraph{Integrals.}
With $C$ the (canonicalized, Sec.~\ref{sec:canon}) MO coefficients,
$h^{\rm MO}=C^{T}h^{\rm AO}C$ and the chemist-notation two-electron integrals
$(p_sr_s|q_ss_s)$ from \code{ao2mo.kernel} restored to a four-index array,
\begin{align}
  h_{pq}&=h^{\rm MO}_{p_sq_s}\,\delta_{\sigma_p\sigma_q},
  \label{eq:h-so}\\
  V[p,r,q,s]\equiv v_{pr,qs}&=(p_sr_s|q_ss_s)\,\delta_{\sigma_p\sigma_r}\delta_{\sigma_q\sigma_s}
  =\braket{pq|rs},
  \label{eq:V-so}\\
  \bar v_{pr,qs}&=v_{pr,qs}-v_{ps,qr}
  \qquad(\code{V - V.transpose(0,3,2,1)}),
  \label{eq:vbar}
\end{align}
so that the interaction is $\tfrac12\sum v_{pr,qs}\,a_p^\dagger a_q^\dagger a_sa_r
=\tfrac14\sum\bar v_{pr,qs}\,a_p^\dagger a_q^\dagger a_sa_r$ (WFL
convention).  The residual interaction removes the all-active elements,
\begin{equation}
  \vR_{pr,qs}=\bigl(1-\delta_{p\in\mathcal A}\delta_{r\in\mathcal A}\delta_{q\in\mathcal A}\delta_{s\in\mathcal A}\bigr)\,v_{pr,qs},
  \qquad
  \vbR_{pr,qs}=\vR_{pr,qs}-\vR_{ps,qr}.
  \label{eq:vR}
\end{equation}

\paragraph{CI vectors and operators.}
CI vectors of a sector with $(n_\alpha,n_\beta)$ electrons in $n$ spatial
orbitals are PySCF arrays of shape $(N_{n_\alpha},N_{n_\beta})$,
$N_k=\binom{n}{k}$ (\code{cistring.num\_strings}).  The elementary
operations are
\begin{align}
  \hat a_p^\dagger\ket{v}&=\code{cre\_a}(v,n,(n_\alpha,n_\beta),p_s)\ \text{if }\sigma_p=0,\qquad
  \code{cre\_b}(v,n,(n_\alpha,n_\beta),p_s)\ \text{if }\sigma_p=1,
  \label{eq:cre}\\
  \hat a_p\ket{v}&=\code{des\_a}(v,n,(n_\alpha,n_\beta),p_s)\ \text{if }\sigma_p=0,\qquad
  \code{des\_b}(v,n,(n_\alpha,n_\beta),p_s)\ \text{if }\sigma_p=1,
  \label{eq:des}
\end{align}
with target sector $(n_\alpha{\pm}1,n_\beta)$ for $\sigma_p=0$ and
$(n_\alpha,n_\beta{\pm}1)$ for $\sigma_p=1$, returning \code{None} when the
target sector is empty ($n_\alpha=n$ for a
creation, $n_\alpha=0$ for an annihilation).  These routines are adjoint
pairs and satisfy the fermionic anticommutation relations across spins
(verified numerically before use).  The Hamiltonian $H(h_1,\mathrm{eri})$ of
a sector is applied as
\begin{equation}
  H\ket{v}=\code{contract\_2e}\bigl(\code{absorb\_h1e}(h_1,\mathrm{eri},n,(n_\alpha,n_\beta),\tfrac12),\,v\bigr),
  \qquad H\ket{v}=0\ \text{if }n_\alpha+n_\beta=0,
  \label{eq:applyH}
\end{equation}
and the dense Hamiltonian matrix of a sector is obtained by applying
Eq.~\eqref{eq:applyH} to every unit vector and symmetrizing,
$H\leftarrow\tfrac12(H+H^{T})$.  Overlaps $\braket{v|w}$ are
\code{numpy.vdot} of the flattened arrays.

\section{Dyall reference (\code{dyall.py})}
\label{sec:dyall}

\subsection{Molecule and CAS calculation}
\begin{itemize}[leftmargin=2em]
\item Linear H$_4$: atoms at $(0,0,k\,r)$, $k=0,\ldots,3$, $r=1.0$~\AA{}
  (\code{build\_h4}); RHF followed by CASSCF$(n_{\rm el},n_{\rm cas})$
  (\code{run\_casscf}).
\item Ozone: $R(\mathrm{O{-}O})=1.278$~\AA{} and $\theta=116.8^\circ$; the
  central O is at the origin and the terminal O atoms at
  $(0,\pm R\sin\frac{\theta}{2},-R\cos\frac{\theta}{2})$; point-group
  symmetry is switched on (\code{build\_o3}).  RHF is followed by CASCI in
  the RHF orbitals with PySCF's default choice of the $n_{\rm cas}$ orbitals
  starting at $n_{\rm core}=(N-n_{\rm el})/2$ (\code{run\_casci}).
\end{itemize}

\subsection{Dyall Fock operator and canonicalization}
\label{sec:canon}

With the spatial CAS one-particle density $\gamma^{\rm cas}$
(\code{mc.fcisolver.make\_rdm1}) the spatial MO density is
$\gamma^{\rm MO}=\operatorname{diag}(2\cdot\mathbf 1_{\rm core},\gamma^{\rm cas},0)$,
and
\begin{equation}
  F^{\rm AO}=h^{\rm AO}+J[\gamma^{\rm AO}]-\tfrac12K[\gamma^{\rm AO}],\qquad
  \gamma^{\rm AO}=C\gamma^{\rm MO}C^{T},\qquad
  F^{\rm MO}=C^{T}F^{\rm AO}C ,
  \label{eq:fock}
\end{equation}
where $J-\tfrac12K$ is \code{mf.get\_veff(mol,\,dm)} of the RHF object.  In
spin orbitals this is $F_{PQ}=h_{PQ}+\sum_{rs}\bar v_{Pr,Qs}\gamma_{rs}$
with the Dyall density of Eq.~\eqref{eq:gamma-so}.  The core--core and
virtual--virtual blocks of $F^{\rm MO}$ are diagonalized, each irrep
separately when orbital symmetry labels are available (so that the labels
survive), the eigenvectors are sorted by eigenvalue, and the core and virtual
columns of $C$ are rotated accordingly; active orbitals are not changed.
$F^{\rm MO}$, $h^{\rm MO}$ and the ERIs are then recomputed in the rotated
basis, and
\begin{equation}
  \epsilon_p=F^{\rm MO}_{p_sp_s}\qquad(\text{used for }p\in\mathcal I).
  \label{eq:eps}
\end{equation}

\subsection{Dyall partition}

The active Dyall Hamiltonian $\hat H^{\act}_D$ has the spatial one-body part
(core index $k$)
\begin{equation}
  h^{\rm eff}_{xy}=h^{\rm MO}_{xy}+\sum_k\bigl[2(xy|kk)-(xk|ky)\bigr]
  \label{eq:heff}
\end{equation}
(checked against \code{mc.get\_h1eff}) and the all-active ERIs
$(xy|zw)$.  The full-space one- and two-body parts of $\hat H_D$, used only
for exact checks, are
\begin{equation}
  h^{D}_{p_sq_s}=\begin{cases}\epsilon_{p}\delta_{p_sq_s}&p_s,q_s\ \text{inactive}\\
  h^{\rm eff}_{p_sq_s}&p_s,q_s\ \text{active}\\ 0&\text{otherwise}\end{cases},
  \qquad
  \mathrm{eri}^{D}_{p_sr_sq_ss_s}=\begin{cases}(p_sr_s|q_ss_s)&\text{all four active}\\0&\text{otherwise}\end{cases}.
  \label{eq:HD}
\end{equation}
The one-body residual interaction in spin orbitals is
\begin{equation}
  u_{pq}=h_{pq}-h^{D}_{pq}\qquad(h^{D}\ \text{expanded to spin orbitals as in Eq.~\eqref{eq:h-so}}),
  \label{eq:u}
\end{equation}
which reproduces WFL Eqs.~(16)--(17) because the inactive blocks of
$F^{\rm MO}$ are diagonal after canonicalization.

\subsection{Active-space sectors and charged poles}
\label{sec:poles}

With $(n_a,n_b)$ the active electron numbers, the sectors
\begin{align}
  \mathcal S_0&=(n_a,n_b),\qquad
  \mathcal S_+=\{(n_a{+}1,n_b),(n_a,n_b{+}1)\}\cap\{n_\sigma\le n_{\rm cas}\},
  \label{eq:sectors}\\
  \mathcal S_-&=\{(n_a{-}1,n_b),(n_a,n_b{-}1)\}\cap\{n_\sigma\ge0\}\nonumber
\end{align}
are diagonalized densely with $H(h^{\rm eff},\mathrm{eri}^{\rm cas})$:
$H^{(\mathcal S)}U^{(\mathcal S)}=U^{(\mathcal S)}\operatorname{diag}(E^{(\mathcal S)})$.
The reference is $E_0=E^{(\mathcal S_0)}_0$, $\ket{\Xi_0}=U^{(\mathcal S_0)}_{:,0}$
(its overlap with \code{mc.ci} is reported).  Precomputed vectors:
\begin{equation}
  \ket{c_x}=\hat a_x^\dagger\ket{\Xi_0},\qquad \ket{d_x}=\hat a_x\ket{\Xi_0},\qquad x\in\{0,\ldots,n_{\act}-1\},
  \label{eq:cd}
\end{equation}
each with its sector (or \code{None}).  The charged poles are, for every
sector $\mathcal S\in\mathcal S_\pm$ and every eigenstate $k$ with vector
$\ket{\Xi_k}$,
\begin{align}
  \text{attachment }(\text{sign }+1):\quad&\kappa_\alpha=E^{(\mathcal S)}_k-E_0,&
  [d_\alpha]_x&=\braket{c_x|\Xi_k}=\braket{\Xi_0|\hat a_x|\Xi_k},
  \label{eq:pole-att}\\
  \text{removal }(\text{sign }-1):\quad&\kappa_\alpha=-(E^{(\mathcal S)}_k-E_0),&
  [d_\alpha]_x&=\braket{\Xi_k|d_x}=\braket{\Xi_k|\hat a_x|\Xi_0},
  \label{eq:pole-rem}
\end{align}
where $[d_\alpha]_x$ is set to zero when the sector of $c_x$ (respectively
$d_x$) differs from $\mathcal S$.  Each pole $\alpha$ stores (sign, sector, $\kappa_\alpha$, $\ket{\Xi_k}$,
$d_\alpha$); the matrix $D_{\act}=[d_\alpha]$ has shape $n_{\act}\times n_{\rm poles}$.
The active sum rule $D_{\act}D_{\act}^{T}=\mathbf 1$ follows from completeness
of the $N_a\pm1$ eigenstates and is checked.

\subsection{Density, static self-energy, transition tensors}

\begin{align}
  \gamma^{\act}_{xy}&=\braket{d_x|d_y}=\braket{\Xi_0|\hat a_x^\dagger\hat a_y|\Xi_0}
  \ (\text{sectors equal, else }0),\qquad
  \gamma_{pq}=\begin{cases}\delta_{pq}&p,q\in\mathcal C\\ \gamma^{\act}_{x_\ell y_\ell}&p,q\in\mathcal A\\0&\text{otherwise,}\end{cases}
  \label{eq:gamma-so}\\
  [\Sigma_{\rm stat}]_{pq}&=u_{pq}+\sum_{rs}\vbR_{pq,rs}\,\gamma_{rs}
  \qquad(\text{WFL Eq.~(20), }\Sigma_1^{11}),
  \label{eq:sigma-stat}\\
  C_3[\alpha,z,y,w]&=\begin{cases}
  \braket{\hat a_z^\dagger\hat a_w^\dagger\hat a_y\Xi_0\,|\,\Xi_\alpha}
  =\braket{\Xi_0|\hat a_y^\dagger\hat a_w\hat a_z|\Xi_\alpha}&\text{attachment poles}\\[1mm]
  \braket{\Xi_\alpha\,|\,\hat a_y^\dagger\hat a_w\hat a_z\Xi_0}&\text{removal poles}
  \end{cases}
  \label{eq:C3}\\
  \mathcal C_{\alpha;z,yw}&=\tfrac12\,C_3[\alpha,z,y,w]-\gamma^{\act}_{yw}\,[d_\alpha]_z ,
  \label{eq:Cc}
\end{align}
where the operator strings are applied right to left with
Eqs.~\eqref{eq:cre}--\eqref{eq:des} and contribute only when the resulting
sector equals the pole's sector.  Equation~\eqref{eq:Cc} is the connected
tensor $\mathcal C^{D,\pm}_{\mu;z,yw}$ of the Dyall--Hedin notes.

\subsection{Neutral transition densities}

For the excited states $\mu\ge1$ of $\mathcal S_0$,
\begin{equation}
  \omega_\mu=E^{(\mathcal S_0)}_\mu-E_0,\qquad
  \rho^\mu_{xy}=\braket{\Xi_\mu|\hat a_x^\dagger\hat a_y|\Xi_0}
  =\braket{\Xi_\mu|\hat a_x^\dagger d_y}\quad(\sigma_x=\sigma_y,\ \text{else }0).
  \label{eq:tdm}
\end{equation}

\subsection{Pole representation of $G_D$}

\begin{align}
  T_D&=\bigl[\,e_P\ (P\in\mathcal I)\ \big|\ d_\alpha\ (\text{in the active rows})\,\bigr]\in\mathbb R^{n_{\so}\times(n_I+n_{\rm poles})},
  \label{eq:GD-pole}\\
  K_D&=(\epsilon_P,\ \kappa_\alpha),\qquad
  G_D(z)=T_D\,(z-K_D)^{-1}T_D^{T},\nonumber
\end{align}
with $n_I=|\mathcal I|$.  $T_DT_D^{T}=\mathbf 1$ (zeroth moment) is checked.

\section{MR-RPA screening and the MR-$GW$ bath (\code{mrrpa.py})}
\label{sec:mrrpa}

\subsection{Channels of the reference polarizability}

Each channel $c$ has an energy $\omega_c$ and a transition-density matrix
$\rho^c\in\mathbb R^{n_{\so}\times n_{\so}}$, $\rho^c_{pr}=\braket{\Phi_c|\hat p^\dagger\hat r|\Phi_0}$:
\begin{align}
  \text{(cv)}&\ i\in\mathcal C,\ a\in\mathcal V,\ \sigma_a=\sigma_i: & \rho_{ai}&=1, & \omega&=\epsilon_a-\epsilon_i;\\
  \text{(ca)}&\ i\in\mathcal C,\ \alpha\ \text{attachment}: & \rho_{x i}&=[d_\alpha]_{x_\ell}\ (x\in\mathcal A), & \omega&=\kappa_\alpha-\epsilon_i;\\
  \text{(av)}&\ a\in\mathcal V,\ \alpha\ \text{removal}: & \rho_{a x}&=[d_\alpha]_{x_\ell}, & \omega&=\epsilon_a-\kappa_\alpha;\\
  \text{(aa)}&\ \mu\ge1: & \rho_{xy}&=\rho^\mu_{x_\ell y_\ell}, & \omega&=\omega_\mu .
  \label{eq:channels}
\end{align}
Channels with $\|\rho^c\|_F<10^{-10}$ are discarded (this removes the
spin-forbidden ca/av combinations).  The (cv) channel appears once per
spin-conserving pair; the ca/av channels contain all attachment/removal poles
of both $S_z$ sectors and prune themselves through the norm test.

\subsection{Direct MR-RPA equations}

With $K$ channels, $\rho_f\in\mathbb R^{K\times n_{\so}^2}$ the flattened
densities and $\vR_m\in\mathbb R^{n_{\so}^2\times n_{\so}^2}$ the residual
interaction with row pair $(p,r)$ and column pair $(q,s)$,
\begin{align}
  [V\rho]^{c}_{pr}&=\sum_{qs}\vR_{pr,qs}\,\rho^{c}_{qs},\qquad
  [V\rho^{T}]^{c}_{pr}=\sum_{qs}\vR_{pr,qs}\,\rho^{c}_{sq},\\
  B_{cc'}&=\sum_{pr}\rho^{c}_{pr}\,[V\rho]^{c'}_{pr},\qquad
  A_{cc'}=\omega_c\delta_{cc'}+\sum_{pr}\rho^{c}_{pr}\,[V\rho^{T}]^{c'}_{pr}
  \qquad(\text{WFL Eqs.~(S30)--(S31)}),
  \label{eq:AB}
\end{align}
followed by $A\leftarrow\tfrac12(A+A^{T})$, $B\leftarrow\tfrac12(B+B^{T})$
(the asymmetry before symmetrization is reported).  The RPA problem
$\bigl(\begin{smallmatrix}A&B\\B&A\end{smallmatrix}\bigr)\bigl(\begin{smallmatrix}X\\Y\end{smallmatrix}\bigr)
=\Omega\bigl(\begin{smallmatrix}X\\-Y\end{smallmatrix}\bigr)$ is solved as
\begin{align}
  A-B&=LL^{T}\ (\text{Cholesky}),\qquad
  L^{T}(A+B)L\,Z=Z\,\Omega^{2},
  \label{eq:rpa}\\
  R&\equiv X+Y=LZ\,\Omega^{-1/2},\qquad
  S\equiv X-Y=(A+B)R\,\Omega^{-1},\qquad R^{T}S=\mathbf 1 ,\nonumber
\end{align}
with $\Omega>0$ and $Z^{T}Z=\mathbf 1$.  If $A-B$ is not positive definite
the general form is used instead: eigenpairs of
$\bigl(\begin{smallmatrix}A&B\\-B&-A\end{smallmatrix}\bigr)$ with
$\operatorname{Re}\Omega>0$, $|\operatorname{Im}\Omega|<10^{-8}\max(1,|\Omega|)$
and positive norm $X^{T}X-Y^{T}Y>0$ are kept, sorted, and normalized to
unit norm; a warning is issued if fewer than $K$ such solutions exist.

\subsection{Screened-interaction couplings and static $W$}

\begin{align}
  \bar\rho_{I,qs}&=\sum_c\rho^{c}_{qs}\,R_{cI},\qquad
  M_{pr,I}=\sum_{qs}\vR_{pr,qs}\,\bar\rho_{I,qs}\qquad(\text{WFL Eq.~(S43)}),
  \label{eq:M}\\
  w^{c}_{pr,qs}(0)&=-2\sum_I\frac{M_{pr,I}M_{qs,I}}{\Omega_I},\qquad
  W^{0}_{pr,qs}=\vR_{pr,qs}+w^{c}_{pr,qs}(0)\qquad(W_D(\omega=0)),
  \label{eq:W0}
\end{align}
and the mode norms $m_I=\bigl(\sum_{pr}M_{pr,I}^2\bigr)^{1/2}$.  The
frequency-dependent correlation part, used only in checks, is
$w^{c}_{pr,qs}(z)=\sum_IM_{pr,I}M_{qs,I}\bigl[(z-\Omega_I)^{-1}-(z+\Omega_I)^{-1}\bigr]$.

\subsection{The MR-$GW$ pole bath}
\label{sec:bath}

Modes with $m_I\le10^{-12}\max_Jm_J$ are dropped (they are decoupled from the
residual interaction, e.g.\ spin-flip modes).  For the remaining modes the
retarded correlation self-energy of WFL Eq.~(S52),
$\Sigma^{c}(z)=A_B\,(z-E_B)^{-1}A_B^{T}$, is built from the columns
\begin{equation}
\begin{array}{lll}
  \Lambda=(a,I),\ a\in\mathcal V: & E_\Lambda=\epsilon_a+\Omega_I, & [A_B]_{p\Lambda}=M_{ap,I},\\[0.5mm]
  \Lambda=(\alpha{+},I): & E_\Lambda=\kappa_\alpha+\Omega_I, & [A_B]_{p\Lambda}=\sum_{x\in\mathcal A}[d_\alpha]_{x_\ell}\,M_{xp,I},\\[0.5mm]
  \Lambda=(i,I),\ i\in\mathcal C: & E_\Lambda=\epsilon_i-\Omega_I, & [A_B]_{p\Lambda}=M_{pi,I},\\[0.5mm]
  \Lambda=(\alpha{-},I): & E_\Lambda=\kappa_\alpha-\Omega_I, & [A_B]_{p\Lambda}=\sum_{x\in\mathcal A}[d_\alpha]_{x_\ell}\,M_{px,I},
\end{array}
\label{eq:bath}
\end{equation}
and columns with $\|[A_B]_{:\Lambda}\|\le10^{-12}\max_{\Lambda'}\|[A_B]_{:\Lambda'}\|$
are dropped.  The full retarded MR-$GW$ self-energy is
\begin{equation}
  \Sigma^{\rm MR\text{-}GW}(z)=\Sigma_{\rm stat}+A_B\,(z-E_B)^{-1}A_B^{T}.
  \label{eq:sigma-mrgw}
\end{equation}

\section{Hermitian pole-space matrix (\code{hermitian.py})}
\label{sec:hermitian}

\subsection{Top space}

With $T_D$ and $K_D$ from Eq.~\eqref{eq:GD-pole}, $n_{\rm top}=n_I+n_{\rm poles}$,
\begin{align}
  K&=\operatorname{diag}(K_D)+T_D^{T}\Sigma_{\rm stat}T_D
  +\begin{pmatrix}0&K_B\\K_B^{T}&0\end{pmatrix},\qquad K\leftarrow\tfrac12(K+K^{T}),
  \label{eq:K}\\
  [K_B]_{P\alpha}&=\sum_{z,y,w\in\mathcal A}\mathcal I_{Pz,yw}\,\mathcal C_{\alpha;z,yw}
  \qquad(P\in\mathcal I),
  \label{eq:KB}\\
  \mathcal I_{Pz,yw}&=\begin{cases}
  0&\code{mixed='none'}\\
  \vbR_{Pz,yw}&\code{mixed='bare'}\\
  \vR_{Pz,yw}-W^{0}_{Pw,yz}&\code{mixed='screened'}\quad(\mathcal I^{\rm stat}_D=v_R-[W_D(0)]^{x}).
  \end{cases}
  \label{eq:kernel}
\end{align}
The pair $(P,\alpha)$ is ``same-sector'' if ($P\in\mathcal V$ and $\alpha$ is
an attachment pole) or ($P\in\mathcal C$ and $\alpha$ is a removal pole).
With \code{cross\_sector=False} the inactive--active blocks of both
$T_D^{T}\Sigma_{\rm stat}T_D$ and $K_B$ are multiplied by the same-sector
mask; the default keeps all pairs (Objective~II).  \code{include\_static=False}
drops $T_D^{T}\Sigma_{\rm stat}T_D$.

\subsection{Full matrix, resolvent and spectral function}

With the bath of Eq.~\eqref{eq:bath}, $\tilde A=T_D^{T}A_B\in\mathbb R^{n_{\rm top}\times n_B}$,
\begin{equation}
  \bbH=\begin{pmatrix}K&\tilde A\\ \tilde A^{T}&\operatorname{diag}(E_B)\end{pmatrix},\qquad
  \bbT=(T_D\ \ 0)\in\mathbb R^{n_{\so}\times(n_{\rm top}+n_B)},
  \label{eq:Hfull}
\end{equation}
applied matrix-free as
$(\bbH v)_{\rm top}=Kv_{\rm top}+\tilde Av_B$, $(\bbH v)_B=\tilde A^{T}v_{\rm top}+E_B\circ v_B$,
with diagonal $(\operatorname{diag}K,\ E_B)$.  Without bath ($n_B=0$) the
matrix is $K$.  Dense diagonalization $\bbH U=U\operatorname{diag}(E)$ gives
\begin{equation}
  Z_\lambda=\bbT U_{:,\lambda},\qquad w_\lambda=\|Z_\lambda\|^2,\qquad
  G(z)=\sum_\lambda\frac{Z_\lambda Z_\lambda^{T}}{z-E_\lambda},\qquad
  A(\omega)=\sum_\lambda Z_\lambda Z_\lambda^{T}\,\frac{\eta/\pi}{(\omega-E_\lambda)^2+\eta^2}
  \label{eq:spectral}
\end{equation}
(\code{poles}, \code{spectral\_function}; the trace is returned unless the
full matrix is requested).  Without diagonalization (large baths),
\begin{equation}
  G(z)=T_D\Bigl[z-K-\tilde A\,(z-E_B)^{-1}\tilde A^{T}\Bigr]^{-1}T_D^{T},\qquad
  A(\omega)=-\tfrac1\pi\operatorname{Im}\operatorname{Tr}G(\omega+i\eta)
  \label{eq:schur}
\end{equation}
(\code{greens\_function}, \code{spectral\_function\_resolvent}).  Spectral
moments: $M_0=T_DT_D^{T}$, $M_1=T_DKT_D^{T}$.

\subsection{Bath with the top-space Hamiltonian in the one-boson sector}
\label{sec:dressed-bath}

With \code{bath\_hamiltonian='top'} the bath states are $\ket{n}\otimes\ket{1_I}$
for \emph{all} top states $n$ and all coupled modes $I$ ($m_I>10^{-12}\max_Jm_J$),
the top-space Hamiltonian $K$ of Eq.~\eqref{eq:K} acts inside the one-boson
sector, and the couplings are the boson vertices between top states:
\begin{align}
  B_I&=K+\operatorname{diag}(\sigma_n)\,\Omega_I,\qquad
  [g_I]_{mn}=\sum_{pr}[T_D]_{pm}M_{pr,I}[T_D]_{rn}=t_m^{T}M_It_n,
  \label{eq:dressed-bath}\\
  \bbH&=\begin{pmatrix}K&g_1&g_2&\cdots\\ g_1^{T}&B_1&&\\ g_2^{T}&&B_2&\\ \vdots&&&\ddots\end{pmatrix},\nonumber
\end{align}
where $\sigma_n=+1$ for virtual orbitals and attachment poles and $-1$ for
core orbitals and removal poles.  The diagonal bath of Eq.~\eqref{eq:bath}
is the special case in which $B_I$ is replaced by its diagonal
$\kappa_n+\sigma_n\Omega_I$ (and zero-coupling states are dropped); the two
coincide when $K$ is diagonal.  The bath is applied as
$(B\,v)_{(nI)}=\sum_{n'}K_{nn'}v_{(n'I)}+\sigma_n\Omega_Iv_{(nI)}$, the
diagonal is $K_{nn}+\sigma_n\Omega_I$, and the resolvent uses
\begin{equation}
  G(z)=T_D\Bigl[z-K-\sum_Ig_I\,(z-B_I)^{-1}g_I^{T}\Bigr]^{-1}T_D^{T}.
  \label{eq:dressed-schur}
\end{equation}
Without cross-sector couplings ($[K,\sigma]=0$) the bath Schur complement
equals $T_D^{T}\Sigma^{GW}[G_{\rm init}](z)T_D$ with
\begin{equation}
  \Sigma^{GW}[G_{\rm init}]_{pq}(z)=\sum_{k,I}\frac{[M_IZ_k]_p[M_IZ_k]_q}{z-\lambda_k-\sigma_k\Omega_I},
  \qquad K=U\operatorname{diag}(\lambda)U^{T},\ Z_k=T_DU_{:,k},\ \sigma_k=U_{:,k}^{T}\sigma U_{:,k},
  \label{eq:dressed-equiv}
\end{equation}
i.e.\ the MR-$GW$ self-energy built with the poles of the Eq.~(7) propagator
of the preliminary section instead of the Dyall poles
(\code{dressed\_bath\_self\_energy}; checked).

\subsection{Guess vectors for root following}

\code{unit\_vector(k)} is $e_k$ in the $n_{\rm top}+n_B$ space;
\code{index\_inactive(P)} and \code{index\_active\_pole($\alpha$)} give the
top-space positions.  \code{orbital\_guess($p$, sector)} is the row
$\bbT_{p,:}$ (the components of orbital $p$ on all poles); for an active $p$
and sector \code{'attach'}/\code{'remove'} only the active poles of that sign
are kept; the vector is normalized.  The two H$_4$ principal targets use the
HOMO and LUMO spatial orbitals $p_s=N/2-1$ and $N/2$ (spin $\alpha$).

\subsection{Hall-form insertion (illustration of the non-PSD Green's function)}

For $z=\omega+i\eta$, with $K_B$ from Eq.~\eqref{eq:KB} (bare kernel by
default),
\begin{align}
  [\Sigma^{12}(z)]_{Px}&=\sum_\alpha\frac{[K_B]_{P\alpha}[d_\alpha]_{x_\ell}}{z-\kappa_\alpha}\ (P\in\mathcal I,\ x\in\mathcal A),\qquad
  \Sigma^{21}=(\Sigma^{12})^{T},\nonumber\\
  \Sigma^{11}(z)&=\Sigma_{\rm stat}+A_B(z-E_B)^{-1}A_B^{T},\qquad
  M(z)=(\mathbf 1-\Sigma^{21})^{-1}\,G_D(z)\,(\mathbf 1-\Sigma^{12})^{-1},\nonumber\\
  G_{\rm Hall}(z)&=\bigl[M(z)^{-1}-\Sigma^{11}(z)\bigr]^{-1}.
  \label{eq:hall}
\end{align}

\section{Davidson solver with root following (\code{davidson.py})}
\label{sec:davidson}

Input: \code{matvec} (block application of a real symmetric $\bbH$ of
dimension $n$), the diagonal $D$, $m$ guess vectors $g_1,\ldots,g_m$
(normalized), tolerance $\tau$ (scripts: $10^{-7}$ or $10^{-8}$), maximum
subspace $n_{\max}=\min(n,\max(12m,40))$, preconditioner floor
$\delta=10^{-3}$.
\begin{enumerate}[leftmargin=2em]
\item $V\leftarrow$ orthonormal columns of $[g_1\cdots g_m]$ (QR, columns
  with $|R_{jj}|\le10^{-10}$ dropped; an error is raised if fewer than $m$
  remain); $W\leftarrow\bbH V$; reference vectors $F\leftarrow[g_j]$.
\item Iterate ($\le200$ times): $H_s=\tfrac12(V^{T}W+W^{T}V)$,
  $H_sS=S\operatorname{diag}(\theta)$; Ritz vectors $X=VS$, $\bbH X=WS$.
\item Selection (\code{mode='overlap'}): for $j=1,\ldots,m$ in order, choose
  the not yet selected Ritz index $l_j=\arg\max_l|F_{:,j}^{T}X_{:,l}|$
  (\code{'energy'}: closest $\theta_l$ to the $j$-th target energy;
  \code{'lowest'}: the first $m$).
\item Residuals $r_j=(\bbH X)_{:,l_j}-\theta_{l_j}X_{:,l_j}$; converged if
  $\|r_j\|<\tau$ for all $j$; record the history.
\item For each unconverged $j$: $\delta_j=-r_j\oslash(D-\theta_{l_j})$, where
  entries of $D-\theta$ with $|D-\theta|<\delta$ are replaced by
  $\pm\delta$ (sign preserved); orthogonalize $\delta_j$ twice against $V$
  and against the already accepted new vectors; accept if its norm exceeds
  $10^{-8}$ of the norm before orthogonalization (and $10^{-14}$) and
  normalize.  If the preconditioned correction lies in the current subspace
  (this happens when the relevant block of $\bbH$ is exactly diagonal, where
  $-r_j\oslash(D-\theta)$ is parallel to the Ritz vector), the residual
  $r_j$ itself is orthogonalized and used instead.
\item If no new vector was accepted, stop.  If $\dim V+n_{\rm new}>n_{\max}$,
  restart with $V\leftarrow$ orthonormal columns of $[X_{:,l_1}\cdots X_{:,l_m}\ \delta\ldots]$
  and $W\leftarrow\bbH V$; otherwise append.  Set
  $F\leftarrow[X_{:,l_j}]$ (root following) and return to step~2.
\end{enumerate}
Output: $\theta_{l_j}$, $X_{:,l_j}$, convergence flags, iteration count,
residual norms.  Spectral weights of the roots are $\|\bbT X_{:,l_j}\|^2$ and
their overlaps with the guesses $|g_j^{T}X_{:,l_j}|$.

\section{Full-CI reference and validation identities (\code{fci\_reference.py})}
\label{sec:fci}

\subsection{Exact Green's function}

For the full orbital space ($n_{\rm mo}$ orbitals, $(N_\alpha,N_\beta)$
electrons, integrals $h_1=h^{\rm MO}$, $\mathrm{eri}=(pr|qs)$, or the
$\hat H_D$ or $\hat H_D+\lambda\hat V_R$ integrals when given) the sectors
$\mathcal S_0,\mathcal S_\pm$ of Eq.~\eqref{eq:sectors} are diagonalized
densely, $\ket{\Psi_0}$ and $E_0$ are taken from $\mathcal S_0$, and for
every spin orbital $p$ and eigenstate $k$ of $\mathcal S_\pm$
\begin{equation}
  \text{attachment: }\ \kappa=E_k-E_0,\ t_p=\braket{\hat a_p^\dagger\Psi_0|\Psi_k};\qquad
  \text{removal: }\ \kappa=-(E_k-E_0),\ t_p=\braket{\Psi_k|\hat a_p\Psi_0},
  \label{eq:fci-poles}
\end{equation}
giving $T\in\mathbb R^{n_{\so}\times n_{\rm poles}}$, $K$, $G(z)=T(z-K)^{-1}T^{T}$,
$A(\omega)$ as in Eq.~\eqref{eq:spectral} with $Z_\lambda\to T_{:,\lambda}$,
and $M_0=TT^{T}$, $M_1=TKT^{T}$.  ``Principal poles'' of a sector are the
poles with the largest $\|t\|^2$ above $10^{-3}$.

\subsection{Embedding of active-space vectors}

An active-space CI vector of sector $(n_a,n_b)$ is mapped to the full space
with $(n_a+n_{\rm core},n_b+n_{\rm core})$ electrons by setting the amplitude
of every determinant pair to that of its active part: for each active
$\alpha$-string $s$ (\code{cistring.make\_strings(range($n_{\rm cas}$),$n_a$)})
the full string is $s_{\rm full}=(2^{n_{\rm core}}-1)\,|\,(s\ll n_{\rm core})$
with address \code{cistring.str2addr($n_{\rm mo}$, $n_a+n_{\rm core}$, $s_{\rm full}$)},
and likewise for $\beta$; virtual orbitals are empty.  Because the closed
core is an even operator string, all embedded matrix elements equal the
active-space ones up to one common sign per pair of sectors.

\subsection{Same-sector coupling check}

With $K$ from Eq.~\eqref{eq:K} for \code{mixed='bare'}, no bath, and
$E_0^{H}=\braket{\Psi_0|\hat H|\Psi_0}$ ($\ket{\Psi_0}$ the embedded
$\ket{\Xi_0}$, $\hat H$ the full Hamiltonian), the following identities are
tested (all in the full determinant space; $\ket{\Psi_\alpha}$ the embedded
active pole state, $\ket{w_a}=\hat a_a^\dagger\ket{\Psi_0}$ for $a\in\mathcal V$,
$\ket{w_i}=\hat a_i\ket{\Psi_0}$ for $i\in\mathcal C$, $\sigma_\alpha=\pm1$
the pole sign, $\sigma_P=+1$ for virtual and $-1$ for core):
\begin{align}
  \braket{\Psi_\alpha|\hat H|\Psi_\beta}-E_0^{H}\delta_{\alpha\beta}&=\sigma_\alpha K_{\alpha\beta}\quad(\text{same sector}),
  \label{eq:chk-aa}\\
  \braket{w_Q|\hat H|w_P}-E_0^{H}\delta_{PQ}&=\sigma_P K_{QP}\quad(\text{same sector}),
  \label{eq:chk-ii}\\
  \sigma_\alpha\braket{w_P|\hat H|\Psi_\alpha}&=s_\alpha K_{P\alpha},\quad s_\alpha=\pm1\ \text{common to all }P\ \text{of the sector},
  \label{eq:chk-ia}
\end{align}
where $s_\alpha=\operatorname{sign}\bigl(\sum_P\braket{w_P|\hat H|\Psi_\alpha}K_{P\alpha}\bigr)$
absorbs the sector-pair phase of the embedding.

\subsection{First-order Green's-function check}

With $h_1(\lambda)=h^{D}+\lambda(h^{\rm MO}-h^{D})$,
$\mathrm{eri}(\lambda)=\mathrm{eri}^{D}+\lambda(\mathrm{eri}-\mathrm{eri}^{D})$,
$\lambda=10^{-3}$, and $K^{(1)}=K-\operatorname{diag}(K_D)$ from
Eq.~\eqref{eq:K} (bare kernel, no bath, all pole pairs),
\begin{equation}
  \frac{G_{\rm FCI}(\lambda;z)-G_{\rm FCI}(-\lambda;z)}{2\lambda}
  \ \stackrel{?}{=}\ T_D\,R_D(z)\,K^{(1)}\,R_D(z)\,T_D^{T},\qquad R_D(z)=(z-K_D)^{-1},
  \label{eq:chk-fd}
\end{equation}
at $z\in\{-1.2,-0.6,-0.2,0.2,0.6,1.2\}+0.05i$ Ha, together with
$G_{\rm FCI}(0;z)=G_D(z)$.

\section{EKT/ERPA reference (\code{ekt\_erpa.py})}
\label{sec:ekt}

The object replaces the poles, $D_{\act}$, $\kappa$, signs, $C_3$,
$\mathcal C$ and the neutral transition densities of the Dyall reference and
delegates everything else to it.

\subsection{Charged manifolds and generalized eigenproblems}

For each sign and each sector $\mathcal S\in\mathcal S_\pm$ the manifold is
the list of CI vectors
\begin{align}
  \text{primary:}&\quad \{\hat a_x\ket{\Xi_0}\}\ (\text{removal}),\qquad \{\hat a_x^\dagger\ket{\Xi_0}\}\ (\text{attachment}),
  \label{eq:man-primary}\\
  \text{extended (removal):}&\quad \text{primary}\ \cup\ \{\hat a_y^\dagger\hat a_w\hat a_z\ket{\Xi_0}\}_{w>z,\ y},
  \label{eq:man-extended}\\
  \text{extended (attachment):}&\quad \text{primary}\ \cup\ \{\hat a_y\hat a_w^\dagger\hat a_z^\dagger\ket{\Xi_0}\}_{w>z,\ y},\nonumber
\end{align}
keeping only vectors that land in $\mathcal S$ and have norm $\ge10^{-12}$
(operators applied right to left).  With $V=[v_1\cdots v_m]$ (flattened
vectors as rows) and $\hat H=H(h^{\rm eff},\mathrm{eri}^{\rm cas})$,
\begin{equation}
  A_{ij}=\braket{v_i|\hat H-E_0|v_j},\qquad S_{ij}=\braket{v_i|v_j},\qquad
  A,S\leftarrow\text{symmetrized},
  \label{eq:ekt-AS}
\end{equation}
so that for the primary removal manifold $S=\gamma^{\act}$ and for the
primary attachment manifold $S=\mathbf 1-\gamma^{\act T}$ (EKT metrics).
Canonical orthogonalization:
\begin{equation}
  S=U\operatorname{diag}(s)U^{T},\quad
  X_o=U_{:,s>\tau}\operatorname{diag}(s_{s>\tau})^{-1/2},\quad
  \tau=\begin{cases}\tau_{\rm occ}=10^{-8}&\text{primary}\\ \tau_{\rm lin}\max_js_j,\ \tau_{\rm lin}=10^{-10}&\text{extended,}\end{cases}
  \label{eq:canon-orth}
\end{equation}
then $X_o^{T}AX_o\,C_t=C_t\operatorname{diag}(e)$, $C=X_oC_t$
($C^{T}SC=\mathbf 1$), and for each $k$
\begin{equation}
  \ket{\Xi^{\rm EKT}_k}=\sum_jC_{jk}\ket{v_j},\qquad
  \kappa_k=\pm e_k,\qquad
  [d_k]_x=\begin{cases}\braket{\Xi^{\rm EKT}_k|\hat a_x\Xi_0}&\text{removal}\\ \braket{\hat a_x^\dagger\Xi_0|\Xi^{\rm EKT}_k}&\text{attachment,}\end{cases}
  \label{eq:ekt-poles}
\end{equation}
(zero when the sector of $\hat a_x^{(\dagger)}\Xi_0$ differs from $\mathcal S$).
The rank per sector and the number of removed directions are reported; when
the number of EKT poles equals the exact one,
$\max|\kappa^{\rm EKT}_{\rm sorted}-\kappa^{\rm exact}_{\rm sorted}|$ is
reported.  $C_3$ and $\mathcal C$ are recomputed from
Eqs.~\eqref{eq:C3}--\eqref{eq:Cc} with the EKT states.  Active moments
$M_k=\sum_\alpha\kappa_\alpha^kd_\alpha d_\alpha^{T}$, $k=0,\ldots,3$, are
compared with the exact ones.

\subsection{ERPA in the active--active channel}

Natural spin orbitals per spin block: $\gamma^{\act}_{\sigma\sigma}=U_\sigma\operatorname{diag}(n_\sigma)U_\sigma^{T}$,
assembled into $U$ and $n$ over the local active spin orbitals ($\tilde p^\dagger=\sum_xU_{x\tilde p}\hat a_x^\dagger$).
Same-spin excitation vectors and their NO transforms:
\begin{equation}
  \ket{u_{xy}}=\hat a_x^\dagger\hat a_y\ket{\Xi_0}\ (\sigma_x=\sigma_y,\ \text{sector }\mathcal S_0),\qquad
  \ket{\tilde u_{pq}}=\sum_{x,y}U_{xp}U_{yq}\ket{u_{xy}},
  \label{eq:erpa-u}
\end{equation}
for the pair manifold $\mathcal P=\{(p,q):p\ne q,\ \sigma_p=\sigma_q,\ |n_q-n_p|>\tau_{\rm occ}\}$.
The equation-of-motion matrices
\begin{align}
  \mathcal M_{(pq),(rs)}&=\braket{\tilde u_{pq}|\hat H-E_0|\tilde u_{rs}}+\braket{\tilde u_{sr}|\hat H-E_0|\tilde u_{qp}},
  \qquad\mathcal M\leftarrow\tfrac12(\mathcal M+\mathcal M^{T}),
  \label{eq:erpa-M}\\
  \mathcal S_{(pq),(rs)}&=\braket{\tilde u_{pq}|\tilde u_{rs}}-\braket{\tilde u_{sr}|\tilde u_{qp}},
  \label{eq:erpa-S}
\end{align}
are formed.  Equation~\eqref{eq:erpa-M} equals the double commutator
$\braket{0|[E_{qp},[\hat H,E_{rs}]]|0}$ and Eq.~\eqref{eq:erpa-S} the
commutator $\braket{0|[E_{qp},E_{rs}]|0}$, which is
$\delta_{pr}\delta_{qs}(n_q-n_p)$ in the natural-orbital basis; the
deviation of $\mathcal S$ from this diagonal form is reported.  With the excitation pairs $\mathcal E=\{(p,q)\in\mathcal P:n_q>n_p\}$,
their conjugates $\bar{\mathcal E}=\{(q,p)\}$, $N_{(pq)}=n_q-n_p$,
$A=\mathcal M_{\mathcal E\mathcal E}$, $B=\mathcal M_{\mathcal E\bar{\mathcal E}}$
(the identities $\mathcal M_{\bar{\mathcal E}\bar{\mathcal E}}=A$,
$\mathcal M_{\bar{\mathcal E}\mathcal E}=B$ are checked),
\begin{align}
  \tilde A&=N^{-1/2}AN^{-1/2},\qquad \tilde B=N^{-1/2}BN^{-1/2},\qquad
  (\Omega,R,S)=\text{RPA solve of Eq.~\eqref{eq:rpa} for }(\tilde A,\tilde B),
  \label{eq:erpa-solve}\\
  X&=N^{-1/2}\tfrac{R+S}{2},\qquad Y=N^{-1/2}\tfrac{R-S}{2},\nonumber
\end{align}
which is the ERPA problem $\mathcal Mc=\omega\mathcal Sc$ with
$c^{T}\mathcal Sc=1$.  Transition densities $\rho^\nu=\mathcal Sc^\nu$ in the
NO basis and their back-transformation:
\begin{equation}
  \tilde\rho^\nu_{pq}=N_{(pq)}X^\nu_{(pq)},\qquad
  \tilde\rho^\nu_{qp}=-N_{(pq)}Y^\nu_{(pq)}\quad((p,q)\in\mathcal E),\qquad
  \rho^\nu_{xy}=\sum_{pq}U_{xp}U_{yq}\tilde\rho^\nu_{pq}=\braket{\nu|\hat a_x^\dagger\hat a_y|0}.
  \label{eq:erpa-rho}
\end{equation}
$(\Omega_\nu,\rho^\nu)$ replace $(\omega_\mu,\rho^\mu)$ of Eq.~\eqref{eq:tdm}
in the (aa) channel of Eq.~\eqref{eq:channels}; the (ca) and (av) channels
use the EKT poles automatically.  Sum rules checked against the exact
transition densities over all same-spin pairs:
\begin{equation}
  \sum_\nu\bigl[\rho^\nu_{pq}\rho^\nu_{rs}-\rho^\nu_{qp}\rho^\nu_{sr}\bigr]
  \ \text{and}\
  \sum_\nu\Omega_\nu\bigl[\rho^\nu_{pq}\rho^\nu_{rs}+\rho^\nu_{qp}\rho^\nu_{sr}\bigr]
  \quad(\text{ERPA})\ =\ (\text{exact}).
  \label{eq:erpa-sumrules}
\end{equation}

\section{Workflows of the scripts}
\label{sec:scripts}

\subsection{\code{run\_checks.py} (H$_4$, STO-6G, CAS(2,2), CASSCF)}

The 43 checks, with the identity tested and the tolerance:
\begin{center}\small
\begin{tabular}{>{\raggedright\arraybackslash}p{9.4cm}>{\raggedright\arraybackslash}p{5.4cm}}
\toprule
identity & tolerance\\
\midrule
Eq.~\eqref{eq:heff} vs \code{mc.get\_h1eff}; $E_0+E_{\rm core}=E_{\rm CASSCF}$; $|\braket{\Xi_0|\code{mc.ci}}|=1$ & $10^{-10}$, $10^{-8}$, $10^{-8}$\\
$D_{\act}D_{\act}^{T}=\mathbf 1$; $T_DT_D^{T}=\mathbf 1$ & $10^{-10}$\\
$\Sigma_{\rm stat}$: active--active, core--core, virtual--virtual blocks zero; Hermitian & $10^{-10}$, $10^{-8}$, $10^{-8}$, $10^{-10}$\\
MR-RPA: $A,B$ symmetric; $R^{T}S=\mathbf 1$; $\Omega>0$; $W^{0}_{pr,qs}=W^{0}_{qs,pr}$ & $10^{-10}$, $10^{-8}$, $0$, $10^{-10}$\\
Eq.~\eqref{eq:schur} with $K_B=0$ $=[G_D^{-1}-\Sigma^{\rm MR\text{-}GW}]^{-1}$ at four $z$ & $10^{-9}$\\
Eqs.~\eqref{eq:chk-aa}--\eqref{eq:chk-ia} & $10^{-9}$\\
$G_{\rm FCI}(\hat H_D)=G_D$; Eq.~\eqref{eq:chk-fd} (relative) & $10^{-8}$, $10^{-5}$\\
screened mixed variant: $M_0=\mathbf 1$; $w_\lambda\ge0$; $M_1$(bare mixed) $=h+\bar v\gamma_D$ & $10^{-10}$, $10^{-12}$, $10^{-9}$\\
Davidson roots $=$ dense max-overlap eigenvalues; residuals & $10^{-7}$\\
dressed bath: equals the diagonal bath for diagonal $K$; equals Eq.~\eqref{eq:dressed-equiv} without cross-sector couplings; $M_0=\mathbf 1$; $w_\lambda\ge0$; Davidson vs dense & $10^{-10}$, $10^{-9}$, $10^{-10}$, $10^{-12}$, $10^{-7}$\\
EKT (primary): $M_0=\mathbf 1$, $M_1$ exact, poles exact (CAS(2,2)) & $10^{-10}$, $10^{-9}$, $10^{-9}$\\
EKT (extended): $M_0\ldots M_3$ exact; poles exact & $10^{-8}$\\
ERPA sum rules, Eq.~\eqref{eq:erpa-sumrules} & $10^{-9}$\\
EKT/ERPA: linearization identity; $M_0=\mathbf 1$; $w_\lambda\ge0$; Davidson vs dense & $10^{-9}$, $10^{-10}$, $10^{-12}$, $10^{-7}$\\
\bottomrule
\end{tabular}
\end{center}

\subsection{\code{example\_h4.py}}

\begin{enumerate}[leftmargin=2em]
\item Build H$_4$ (default STO-6G, $r=1.0$~\AA, CAS(2,2), CASSCF), the Dyall
  reference, the MR-RPA, and the EKT/ERPA references with the primary and the
  extended charged manifold, each with its own MR-RPA.
\item Build eight Hermitian matrices: CAS (\code{mixed='none'}, no bath),
  MR-$GW$ (\code{'none'}, bath), mixed bare, mixed screened, and the
  EKT/ERPA and EKT(2h1p)/ERPA versions of MR-$GW$ and of the screened mixed
  variant; the full-CI reference (unless \code{--no-fci}).
\item For every matrix, Davidson root following from the guesses
  \code{orbital\_guess(HOMO,'remove')}, \code{orbital\_guess(LUMO,'attach')},
  the unit vector of the first core spin orbital (if any) and of the first
  virtual spin orbital (if any); compare with the dense eigenvalue of
  maximal overlap with the same guess; report weights and overlaps.
\item Print the active-space excitation energies (exact and ERPA), the
  charged poles (exact, primary EKT, extended EKT), the manifold ranks, the
  four most intense poles per sector of every variant and of full CI, the
  sum rules $M_0$, minimal weight, and $\max|M_1-M_1^{\rm FCI}|$ over the
  full matrix and over the inactive--active block.
\item Plot $\operatorname{Tr}A(\omega)$ for $\omega\in[-30,25]$~eV with
  $\eta=0.1$~eV (Eq.~\eqref{eq:spectral}), one panel per variant plus the
  Hall-form insertion of Eq.~\eqref{eq:hall}, full CI dashed, to
  \code{h4\_<basis>\_cas<n>\_<m>\_spectral\_functions.png}.
\end{enumerate}

\subsection{\code{check\_ozone.py}}

\begin{enumerate}[leftmargin=2em]
\item Build O$_3$ (6-31G), RHF; print the orbital energies of orbitals
  6--14 with C$_{2v}$ labels (running index per irrep) against WFL Table~S3;
  CASCI with $(n_{\rm el},n_{\rm cas})$ from \code{--cas} (default 6,4;
  8,5 supported; 8,6 refused unless \code{--force} because the estimated
  bath exceeds $3\times10^{6}$ poles); print the three leading CI
  coefficients.
\item Dyall reference, MR-RPA; EKT/ERPA references (primary, extended) and
  their MR-RPAs.
\item Cation states: for each reference, in the sector $(n_a-1,n_b)$, the
  lowest removal pole of each irrep, the irrep of a pole being the XOR of
  the orbital irreps (PySCF ids) occupied in its dominant determinant.
\item For the variants MR-$GW$, mixed bare, mixed screened (exact
  reference), EKT/ERPA MR-$GW$ and its screened mixed version, and the two
  EKT(2h1p)/ERPA versions: Davidson root following from the unit vectors of
  the $^2A_1$, $^2B_2$, $^2A_2$ cation poles of the respective reference;
  IPs $=-\theta$, weights, overlaps.
\item Table of IPs against WFL Table~S5 (KT, $GW$@RHF, CAS, MR-$GW$,
  EOM-CCSD, DMRG, experiment) with the state ordering, and the deviation
  MR-$GW$(this work) $-$ MR-$GW$(WFL).
\item Optional (\code{--plot}): $\operatorname{Tr}A(\omega)$ on
  $[-20,-8]$~eV with $\eta=0.1$~eV from Eq.~\eqref{eq:schur} for the CASCI
  reference and the three exact-reference variants, with the experimental
  (dotted) and DMRG (dashed) IPs as vertical lines.
\end{enumerate}

\section{Equation-to-function index}
\label{sec:index}

\begin{center}\small
\begin{tabular}{>{\raggedright\arraybackslash}p{4.0cm}>{\raggedright\arraybackslash}p{6.5cm}>{\raggedright\arraybackslash}p{4.5cm}}
\toprule
equations & function / method & module\\
\midrule
\eqref{eq:so}--\eqref{eq:vR} & \code{spin\_orbital\_one\_body}, \code{spin\_orbital\_two\_body}, \code{antisymmetrize}, \code{DyallReference.\_\_init\_\_} & \code{dyall.py}\\
\eqref{eq:cre}--\eqref{eq:applyH} & \code{apply\_cre}, \code{apply\_des}, \code{apply\_hamiltonian}, \code{dense\_ci\_hamiltonian} & \code{dyall.py}\\
\eqref{eq:fock}--\eqref{eq:eps} & \code{DyallReference.\_\_init\_\_}, \code{\_canonicalize\_block} & \code{dyall.py}\\
\eqref{eq:heff}--\eqref{eq:u} & \code{DyallReference.\_\_init\_\_} & \code{dyall.py}\\
\eqref{eq:sectors}--\eqref{eq:pole-rem} & \code{DyallReference.\_\_init\_\_}, \code{Sector}, \code{ActivePole} & \code{dyall.py}\\
\eqref{eq:gamma-so}--\eqref{eq:Cc} & \code{DyallReference.\_\_init\_\_}, \code{\_three\_operator\_amplitudes} & \code{dyall.py}\\
\eqref{eq:tdm} & \code{neutral\_transition\_densities} & \code{dyall.py}\\
\eqref{eq:GD-pole} & \code{pole\_representation}, \code{dyall\_greens\_function} & \code{dyall.py}\\
\eqref{eq:channels} & \code{MRRPA.\_build\_channels} & \code{mrrpa.py}\\
\eqref{eq:AB}--\eqref{eq:rpa} & \code{MRRPA.\_solve}, \code{solve\_rpa}, \code{\_solve\_rpa\_general} & \code{mrrpa.py}\\
\eqref{eq:M}--\eqref{eq:W0} & \code{MRRPA.\_solve} & \code{mrrpa.py}\\
\eqref{eq:bath}--\eqref{eq:sigma-mrgw} & \code{MRRPA.bath}, \code{MRRPA.self\_energy} & \code{mrrpa.py}\\
\eqref{eq:K}--\eqref{eq:kernel} & \code{HermitianGF.\_\_init\_\_}, \code{mixed\_kernel} & \code{hermitian.py}\\
\eqref{eq:Hfull}--\eqref{eq:schur} & \code{matvec}, \code{diagonal}, \code{dense}, \code{eig}, \code{poles}, \code{greens\_function}, \code{spectral\_function(\_resolvent)}, \code{moments}, \code{dense\_roots} & \code{hermitian.py}\\
\eqref{eq:dressed-bath}--\eqref{eq:dressed-equiv} & \code{HermitianGF.\_\_init\_\_} (\code{bath\_hamiltonian='top'}), \code{\_bath\_apply}, \code{\_bath\_block}, \code{dressed\_bath\_self\_energy} & \code{hermitian.py}\\
\eqref{eq:hall} & \code{hall\_insertion\_gf} & \code{hermitian.py}\\
Sec.~\ref{sec:davidson} & \code{davidson\_root\_following} & \code{davidson.py}\\
\eqref{eq:fci-poles} & \code{FCIReference} & \code{fci\_reference.py}\\
embedding, \eqref{eq:chk-aa}--\eqref{eq:chk-fd} & \code{embed\_active\_vector}, \code{check\_same\_sector\_couplings}, \code{check\_first\_order\_gf} & \code{fci\_reference.py}\\
\eqref{eq:man-primary}--\eqref{eq:ekt-poles} & \code{EKTERPAReference.\_build\_ekt} & \code{ekt\_erpa.py}\\
\eqref{eq:erpa-u}--\eqref{eq:erpa-sumrules} & \code{EKTERPAReference.\_build\_erpa}, \code{response\_sum\_rules}, \code{active\_moments} & \code{ekt\_erpa.py}\\
\bottomrule
\end{tabular}
\end{center}

\end{document}
